\documentclass[11pt]{article}

\usepackage[utf8]{inputenc}
\usepackage[T1]{fontenc}
\usepackage{geometry}
\geometry{margin=1in}
\usepackage{amsmath,amssymb}
\usepackage{booktabs}
\usepackage{graphicx}
\usepackage{hyperref}

\title{Empirical Validation Note for Matching GFM}
\author{Working empirical note}
\date{\today}

\begin{document}
\maketitle

\paragraph{Scope.}
This note summarizes the current executable experiment stack for the
matching-GFM project. It is intentionally narrower than the full proposal.
The present implementation covers a synthetic temporal bipartite market with
typed edges, preference recovery, one-to-one matching, and stability-aware
evaluation.

\paragraph{Implemented components.}
The current runner includes:
\begin{itemize}
  \item a synthetic temporal bipartite multigraph generator;
  \item a pointwise gradient-boosted baseline;
  \item a compact graph-matching model with recency weighting, Sinkhorn
  relaxation, and blocking-pair regularization;
  \item matching metrics covering preference recovery, stability, welfare, and
  held-out ranking quality.
\end{itemize}

\IfFileExists{generated_results.tex}{
  \input{generated_results.tex}
}{
  \paragraph{Missing generated results.}
  Run \texttt{python3 experiments/render\_empirical\_note.py} from the project
  root to materialize the current tables and figure references.
}

\paragraph{Interpretive caution.}
These results should be read as proof of viability for the paper scaffold, not
as the final empirical claim of a graph foundation model. The compact matcher is
still a low-parameter approximation to the fuller temporal-graph agenda.

\end{document}
